\centerline{\bf \Huge Банки. База(А+Д)}
\centerline{\bf \Huge }

\begin{flushright}
        \scriptsize{\textbf{tt}}
\end{flushright}

\problem{\textbf{1}} Для покупки квартиры Алексею не хватало 1 209 600 рублей, поэтому в январе 2015 года
он решил взять в банке кредит под 10\% годовых на 2 года. 

Условия пользования кредитом таковы:

— раз в год 15 декабря банк начисляет на оставшуюся сумму долга проценты, то есть долг
увеличивается на 10\%;

— в период с 16 по 31 декабря Алексей обязан перевести в банк некоторый платеж в $x$
рублей.

Чему должен быть равен $x$, чтобы Алексей выплатил долг равными платежами?

\problem{\textbf{2}} В воскресенье, 20 октября, на Пагоде семь гопников остановили Анджура Аркадьевича. После недлительных переговоров АА предложил гопникам интересную схему заёма денежных средств. Каждый понедельник долг увеличивается на 20\%, и со вторника по субботу необходимо выплатить одним платежом часть долга. Гопники три раза переводили одну и ту же сумму на оффшорный счет АА и полностью погасили свой долг. Сколько рублей они изначально занимали у АА, если им пришлось вернуть на 77200 рублей больше суммы, взятой в долг?

\problem{\textbf{3}} Бизнесмен Олег в январе 2016 года взял кредит в банке под 20 \% годовых, причем выплачивать кредит он должен равными суммами в течение трех лет. Сколько рублей в итоге
выплатил Олег банку, если известно, что его переплата по кредиту составила 675 500 рублей?

\problem{\textbf{4}} Руслан взял кредит в банке под $y$ \% годовых. Выплачивать кредит он должен в течение
двух лет равными ежегодными платежами, переводимыми в банк после начисления процентов.
Под какой процент $y$ был взят кредит, если ежегодный платеж составил $\dfrac{81}{136}$ от суммы кредита?

\problem{\textbf{5}} 10 лет назад Григорий брал в банке кредит на 4 года, причем Григорий помнит, что
выплачивал он кредит дифференцированными платежами и переплата по кредиту составила
32,5\% от кредита. Под какой годовой процент был взят тогда кредит?

\problem{\textbf{6}} Семья взяла в банке ипотечный кредит под 10\% годовых на 8 лет. Условия погашения
кредита следующие: по истечении каждого года заемщик погашает банку начисленные проценты за год и $\dfrac{1}{8}$ часть основной суммы. Какую сумму семья взяла в банке, если последний
платеж, которым она полностью погасила кредит, составил 605 тысяч рублей? Ответ дайте в
миллионах рублей.

\problem{\textbf{7}} Банк выдает кредит на следующих условиях:

— раз в год банк начисляет на текущий долг некоторый процент годовых;

— раз в год после начисления процентов клиент обязан внести платеж в счет погашения
кредита, причем платежи вносятся таким образом, чтобы сумма долга уменьшалась каждый
год на одну и ту же величину;

Сколько процентов составит переплата от кредита, если взять такой кредит на 9 лет и отношение наибольшего платежа к наименьшему платежу равно 17 : 9.?

\problem{\textbf{8}} 15 мая планируется взять кредит в банке сроком на 23 месяца. Условия следующие: 

— 1-го числа каждого месяца долг возрастает на $t$\% по сравнению с концом предыдущего
месяца;

— со 2-го по 14-е число каждого месяца необходимо выплатить часть долга;

— 15-го числа каждого месяца долг должен быть на одну и ту же величину меньше долга
на 15-е число предыдущего месяца.

Известно, что общая сумма денег, которую нужно выплатить банку за весь срок кредитования, на 36\% больше, чем сумма, взятая в кредит. Найдите $t$.

\problem{\textbf{9}} В июле планируется взять кредит в банке на сумму 14 млн. рублей на некоторое целое
число лет. Условия его возврата таковы:

– каждый январь долг возрастает на 25\% по сравнению с концом предыдущего года;

– с февраля по июнь каждого года необходимо выплачивать часть долга;

– в июле каждого года долг должен быть на одну и ту же сумму меньше долга на июль
предыдущего года.

На сколько лет взят кредит, если известно, что общая сумма выплат после его погашения
равнялась 24,5 млн. рублей?
